\chapter{1996年上海卷（文）}
\section{单选题}

\begin{problem}
已知集合 \(M=\left\{(x,y)\mid x+y=2\right\}\)，\(N=\left\{(x,y)\mid x-y=4\right\}\)，那么集合 \(M\cap N\) 为（\(\quad\)）

\choices
    {\(x=3\)，\(y=-1\)}
    {\((3,-1)\)}
    {\(\{3,-1\}\)}
    {\(\{(3,-1)\}\)}
\end{problem}

\begin{answer}
D
\end{answer}

\begin{solution}
集合 \(M\cap N\) 中的点必须同时满足两个方程，因此
\[
\begin{cases}
x+y=2,\\
x-y=4.
\end{cases}
\]
两式相加得 \(2x=6\)，所以 \(x=3\)；代入 \(x+y=2\) 得 \(y=-1\). 因而 \(M\cap N\) 只有一个元素，即 \(M\cap N=\left\{(3,-1)\right\}\)，故选 D.
\end{solution}

\begin{problem}
在下列各区间中，函数 \(y=\sin\left(x+\frac{\pi}{4}\right)\) 的单调递增区间是（\(\quad\)）

\choices
    {\(\left[\frac{\pi}{2},\pi\right]\)}
    {\(\left[0,\frac{\pi}{4}\right]\)}
    {\(\left[-\pi,0\right]\)}
    {\(\left[\frac{\pi}{4},\frac{\pi}{2}\right]\)}
\end{problem}

\begin{answer}
B
\end{answer}

\begin{solution}
令 \(t=x+\frac{\pi}{4}\)，则函数为 \(y=\sin t\). 正弦函数在
\(\left[-\frac{\pi}{2}+2k\pi,\frac{\pi}{2}+2k\pi\right]\)
上单调递增. 对选项 A，\(x\in\left[\frac{\pi}{2},\pi\right]\) 时
\(t\in\left[\frac{3\pi}{4},\frac{5\pi}{4}\right]\)，正弦函数在该区间上单调递减；对选项 B，\(t\in\left[\frac{\pi}{4},\frac{\pi}{2}\right]\)，整个区间位于递增区间内；对选项 C，\(t\in\left[-\frac{3\pi}{4},\frac{\pi}{4}\right]\)，跨过 \(t=-\frac{\pi}{2}\)，先减后增；对选项 D，\(t\in\left[\frac{\pi}{2},\frac{3\pi}{4}\right]\)，位于递减区间内. 因此只有选项 B 符合，故选 B.
\end{solution}

\begin{problem}
过点 \((4,0)\) 和点 \((0,3)\) 的直线的倾斜角为（\(\quad\)）

\choices
    {\(\arctan \frac{3}{4}\)}
    {\(\pi-\arctan \frac{3}{4}\)}
    {\(\arctan\left(-\frac{3}{4}\right)\)}
    {\(\pi-\arctan\left(-\frac{3}{4}\right)\)}
\end{problem}

\begin{answer}
B
\end{answer}

\begin{solution}
直线的斜率为
\[
k=\frac{3-0}{0-4}=-\frac{3}{4}
\]
设直线的倾斜角为 \(\theta\)，则 \(0\leq\theta<\pi\)，且 \(\tan\theta=-\frac{3}{4}\). 倾斜角在第二象限，所以
\[
\theta=\pi-\arctan\frac{3}{4}
\]
故选 B.
\end{solution}

\begin{problem}
在下列命题中，真命题是（\(\quad\)）

\choices
    {若直线 \(m\)，\(n\) 都平行于平面 \(\alpha\)，则 \(m\parallel n\)}
    {设 \(\alpha-l-\beta\) 是直二面角. 若直线 \(m\perp l\)，则 \(m\perp\beta\)}
    {若直线 \(m\)，\(n\) 在平面 \(\alpha\) 内的射影依次是一个点和一条直线，且 \(m\perp n\)，则 \(n\) 在 \(\alpha\) 内或 \(n\) 与 \(\alpha\) 平行}
    {设 \(m\)，\(n\) 是异面直线. 若 \(m\) 与平面 \(\alpha\) 平行，则 \(n\) 与 \(\alpha\) 相交}
\end{problem}

\begin{answer}
C
\end{answer}

\begin{solution}
先说明 C 正确. 直线 \(m\) 在平面 \(\alpha\) 内的射影是一个点，所以 \(m\perp\alpha\). 取 \(m\) 的方向为平面 \(\alpha\) 的法向方向. 由于 \(m\perp n\)，直线 \(n\) 的方向在该法向方向上的分量必须为零，因此 \(n\) 平行于平面 \(\alpha\). 于是 \(n\) 若与 \(\alpha\) 相交，就只能全部位于 \(\alpha\) 内；若不相交，则 \(n\parallel\alpha\). 这正是选项 C 的结论.

其余选项都可构造反例. 例如，在平行于 \(\alpha\) 的另一个平面内取两条相交直线 \(m\)，\(n\)，它们都平行于 \(\alpha\)，但不互相平行；在直二面角的两个面分别取为 \(xz\) 平面和 \(yz\) 平面、棱为 \(z\) 轴，取方向向量为 \((1,1,0)\) 的直线 \(m\)，则 \(m\perp l\)，但 \(m\) 不垂直于 \(\beta\)；再取 \(\alpha\) 为 \(z=0\)，取一条位于 \(z=1\) 且平行于 \(x\) 轴的直线 \(m\)，以及一条位于 \(z=2\) 且平行于 \(y\) 轴的直线 \(n\)，则 \(m\)，\(n\) 异面，\(m\parallel\alpha\)，而 \(n\parallel\alpha\) 并不与 \(\alpha\) 相交. 因此真命题是 C.
\end{solution}

\begin{problem}
将椭圆 \(\frac{x^2}{25}+\frac{y^2}{9}=1\) 绕其左焦点按逆时针方向旋转 \(90^\circ\) 后所得的椭圆方程是（\(\quad\)）

\choices
    {\(\frac{(x+4)^2}{25}+\frac{(y-4)^2}{9}=1\)}
    {\(\frac{(x+4)^2}{25}+\frac{(y+4)^2}{9}=1\)}
    {\(\frac{(x+4)^2}{9}+\frac{(y-4)^2}{25}=1\)}
    {\(\frac{(x+4)^2}{9}+\frac{(y+4)^2}{25}=1\)}
\end{problem}

\begin{answer}
C
\end{answer}

\begin{solution}
椭圆的半长轴为 \(5\)，半短轴为 \(3\)，焦半距为 \(c=\sqrt{25-9}=4\)，所以左焦点为 \(F(-4,0)\). 以 \(F\) 为中心逆时针旋转 \(90^\circ\) 时，原中心 \(O(0,0)\) 的相对位置向量 \((4,0)\) 变为 \((0,4)\)，故新中心为 \((-4,4)\). 原来的水平长轴旋转后变为竖直长轴，长半轴仍为 \(5\)，短半轴仍为 \(3\). 因而旋转后的椭圆方程为
\[
\frac{(x+4)^2}{9}+\frac{(y-4)^2}{25}=1
\]
故选 C.
\end{solution}

\begin{problem}
若 \(y=f(x)\) 是定义在 \(\R\) 上的函数，则 \(y=f(x)\) 为奇函数的一个充要条件为（\(\quad\)）

\choices
    {\(f(0)=0\)}
    {对任意的 \(x\in\R\)，\(f(x)=0\) 都成立}
    {存在某一 \(x_0\in\R\)，使得 \(f(x_0)+f(-x_0)=0\)}
    {对任意的 \(x\in\R\)，\(f(x)+f(-x)=0\) 都成立}
\end{problem}

\begin{answer}
D
\end{answer}

\begin{solution}
奇函数的定义是：对定义域内任意 \(x\)，都有 \(f(-x)=-f(x)\). 由于定义域为 \(\R\)，这等价于
\[
f(x)+f(-x)=0
\]
因此“对任意 \(x\in\R\)，\(f(x)+f(-x)=0\)”是奇函数的充要条件. 选项 A 是奇函数的必要条件，但不是充分条件，例如 \(f(x)=x^2\) 满足 \(f(0)=0\) 而不是奇函数. 选项 B 是充分条件，因为零函数是奇函数，但它不是必要条件，例如 \(f(x)=x\) 是奇函数而不满足“对任意 \(x\)，\(f(x)=0\)”. 选项 C 也是必要条件：奇函数取 \(x_0=0\) 即有 \(f(x_0)+f(-x_0)=0\)，但它不是充分条件，例如 \(f(x)=x^2\) 取 \(x_0=0\) 满足该等式而函数仍不是奇函数. 故选 D.
\end{solution}

\begin{problem}
如果 \(\log_a 3>\log_b 3>0\)，那么 \(a\)，\(b\) 间的关系是（\(\quad\)）

\choices
    {\(0<b<a<1\)}
    {\(0<a<b<1\)}
    {\(1<b<a\)}
    {\(1<a<b\)}
\end{problem}

\begin{answer}
D
\end{answer}

\begin{solution}
对数底数必须为正且不等于 \(1\). 由于 \(3>1\)，而 \(\log_a3>0\)，可知 \(a>1\)；同理 \(b>1\). 又
\(\log_a3=\frac{\ln3}{\ln a}\)，\(\log_b3=\frac{\ln3}{\ln b}\).
因为 \(\ln3>0\)，且 \(\ln a\)，\(\ln b\) 均为正，\(\log_a3>\log_b3\) 等价于 \(\frac{1}{\ln a}>\frac{1}{\ln b}\)，从而 \(\ln a<\ln b\)，即 \(a<b\). 因此 \(1<a<b\)，故选 D.
\end{solution}

\begin{problem}
在下列图象中，二次函数 \(y=ax^2+bx\) 与指数函数 \(y=\left(\frac{b}{a}\right)^x\) 的图象只可能是（\(\quad\)）

\choices
    {\choicebitmap[width=0.15\paperwidth]{img/1996/shanghai_liberal/q8_fig1.png}}
    {\choicebitmap[width=0.15\paperwidth]{img/1996/shanghai_liberal/q8_fig2.png}}
    {\choicebitmap[width=0.15\paperwidth]{img/1996/shanghai_liberal/q8_fig3.png}}
    {\choicebitmap[width=0.15\paperwidth]{img/1996/shanghai_liberal/q8_fig4.png}}
\end{problem}

\begin{answer}
A
\end{answer}

\begin{solution}
令 \(q=\frac{b}{a}\). 指数函数的定义要求 \(q>0\)，由图象递减可知 \(0<q<1\). 二次函数可写为
\[
y=ax^2+bx=ax(x+q)
\]
因此抛物线的两个零点为 \(0\) 和 \(-q\)，且 \(-q\in(-1,0)\). 所以图中的抛物线必须与 \(x\) 轴交于原点及 \((-1,0)\) 与 \((0,0)\) 之间的另一点. 第一幅图的抛物线开口向上，另一零点位于该区间，且指数函数递减，满足全部条件；第二幅图的另一零点在正半轴上，第三幅图的另一零点在 \(-1\) 的左侧，第四幅图的另一零点也在正半轴上，均不符合 \(-q\in(-1,0)\). 因此只可能是第一幅图，故选 A.
\end{solution}

\section{填空题}

\begin{problem}
方程 \(\log_{2}\left(9^x-5\right)=\log_{2}\left(3^x-2\right)+2\) 的解是 \(x=\)\fillinblank{}.
\end{problem}

\begin{answer}
\(1\)
\end{answer}

\begin{solution}
令 \(t=3^x\)，则 \(t>0\)，且 \(9^x=t^2\). 对数真数必须为正，所以
\(t^2-5>0\)，\(t-2>0\)，从而 \(t>2\). 由对数的单调性，
\[
t^2-5=4(t-2)
\]
即
\[
t^2-4t+3=0
\]
所以 \(t=1\) 或 \(t=3\). 结合 \(t>2\) 舍去 \(t=1\)，得到 \(3^x=3\)，故 \(x=1\).
\end{solution}

\begin{problem}
函数 \(y=\frac{1}{\sqrt{\log_{\frac{1}{2}}(2-x)}}\) 的定义域是\fillinblank{}.
\end{problem}

\begin{answer}
\(\{x \mid 1 < x < 2\}\)（或 \((1,2)\)）
\end{answer}

\begin{solution}
根式在分母中，必须满足
\(\log_{\frac{1}{2}}(2-x)>0\)，同时对数真数要求 \(2-x>0\). 由于底数 \(\frac{1}{2}\in(0,1)\)，
\(\log_{\frac{1}{2}}u>0\) 等价于 \(0<u<1\). 因而
\[
0<2-x<1
\]
解得 \(1<x<2\). 该条件也保证根式不为零，所以定义域为 \((1,2)\).
\end{solution}

\begin{problem}
方程 \(\sin\left(x-\frac{\pi}{3}\right)=\frac{\sqrt{2}}{2}\) 在 \([0,\pi]\) 上的解是 \(x=\)\fillinblank{}.
\end{problem}

\begin{answer}
\(\frac{7\pi}{12}\)
\end{answer}

\begin{solution}
令 \(t=x-\frac{\pi}{3}\). 当 \(x\in[0,\pi]\) 时，
\(t\in\left[-\frac{\pi}{3},\frac{2\pi}{3}\right]\). 方程
\(\sin t=\frac{\sqrt{2}}{2}\) 的一般解为
\(t=\frac{\pi}{4}+2k\pi\) 或 \(t=\frac{3\pi}{4}+2k\pi\). 在上述区间内只有
\(t=\frac{\pi}{4}\)，所以
\[
x=t+\frac{\pi}{3}=\frac{\pi}{4}+\frac{\pi}{3}=\frac{7\pi}{12}
\]
\end{solution}

\begin{problem}
函数 \(y=x^{-2}\)（\(x<0\)）的反函数是 \(y=\)\fillinblank{}.
\end{problem}

\begin{answer}
\(y=-\sqrt{\frac{1}{x}}\)（\(x>0\)）
\end{answer}

\begin{solution}
原函数的定义域为 \(x<0\)，且 \(y=x^{-2}=\frac{1}{x^2}>0\). 由
\(x^2=\frac{1}{y}\)，并结合原定义域 \(x<0\)，得到
\(x=-\sqrt{\frac{1}{y}}\). 交换 \(x\)，\(y\)，可得反函数
\(y=-\sqrt{\frac{1}{x}}\)，\(x>0\).
\end{solution}

\begin{problem}
与圆 \((x-2)^2+y^2=1\) 外切，且与 \(y\) 轴相切的动圆圆心 \(P\) 的轨迹方程为\fillinblank{}.
\end{problem}

\begin{answer}
\(y^2=6x-3\)
\end{answer}

\begin{solution}
设动圆圆心为 \(P(x,y)\)，由于动圆与 \(y\) 轴相切，动圆半径为 \(r=|x|\). 原圆的圆心为 \(C(2,0)\)，半径为 \(1\). 外切条件给出
\[
\sqrt{(x-2)^2+y^2}=|x|+1
\]
两边均非负，平方不改变等价性，并且
\((x-2)^2+y^2=(|x|+1)^2\)，\(\Longrightarrow\)，\(y^2=4x+2|x|-3\).
若 \(x<0\)，右端为 \(2x-3<0\)，不可能有实数 \(y\). 因而 \(x\geq0\)，于是
\[
y^2=6x-3
\]
反之，若 \(y^2=6x-3\)，则 \(x\geq\frac{1}{2}>0\)，从而

\[
(x-2)^2+y^2=x^2+2x+1=(x+1)^2=(|x|+1)^2
\]

两边均为非负数，故可开平方恢复外切条件. 因此轨迹方程为 \(y^2=6x-3\).
\end{solution}

\begin{problem}
在 \(\left(1+x\right)^2\left(1-x\right)^4\) 的展开式中，\(x^3\) 的系数是（结果用数值表示）\fillinblank{}.
\end{problem}

\begin{answer}
\(4\)
\end{answer}

\begin{solution}
由
\((1+x)^2=1+2x+x^2\)，\((1-x)^4=1-4x+6x^2-4x^3+x^4\).
\(x^3\) 项只能来自第一因子的常数项与第二因子的 \(x^3\) 项、第一因子的 \(2x\) 项与第二因子的 \(x^2\) 项、第一因子的 \(x^2\) 项与第二因子的 \(x\) 项. 因此系数为
\[
-4+2\cdot6-4=4
\]
\end{solution}

\section{选做题：填空题}

\begin{problem}[A]
已知 \(O(0,0)\) 和 \(A(6,3)\) 两点，若点 \(P\) 在直线 \(OA\) 上，且 \(\frac{OP}{PA}=\frac{1}{2}\)，又 \(P\) 是线段 \(OB\) 的中点，则点 \(B\) 的坐标是\fillinblank{}.
\end{problem}

\begin{answer}
\((4,2)\)
\end{answer}

\begin{solution}
设 \(P=tA=(6t,3t)\)，则
\[
\frac{OP}{PA}=\frac{|t|}{|1-t|}=\frac{1}{2}
\]
解得 \(t=\frac{1}{3}\) 或 \(t=-1\). 若按原卷答案所采用的线段 \(OA\) 内分情形，则 \(0<t<1\)，故 \(t=\frac{1}{3}\)，从而 \(P=(2,1)\). 又 \(P\) 是线段 \(OB\) 的中点，所以 \(B=2P=(4,2)\).

严格按题面“直线 \(OA\)”理解时，\(t=-1\) 也满足长度比，此时 \(P=(-6,-3)\)，并得到 \(B=(-12,-6)\). 因此原题文字存在“直线”与标准答案之间的歧义；现答案保留原卷标准答案对应的线段内分解释.
\end{solution}

\begin{problem}[A]
平移坐标轴将抛物线 \(4x^2-8x+y+5=0\) 化为标准方程 \(x^{\prime2}=ay^\prime\)（\(a\neq0\)），则新坐标系的原点在原坐标系中的坐标是\fillinblank{}.
\end{problem}

\begin{answer}
\((1,-1)\)
\end{answer}

\begin{solution}
将方程配方：
\[
4x^2-8x+y+5=4(x-1)^2+y+1=0
\]
即
\[
(x-1)^2=-\frac{1}{4}(y+1)
\]
令新坐标为 \(x'=x-1\)，\(y'=y+1\)，则方程化为
\(x^{\prime2}=-\frac{1}{4}y'\)，符合 \(x^{\prime2}=ay'\) 的形式. 因此新坐标系原点满足
\(x=1\)，\(y=-1\)，其坐标为 \((1,-1)\).
\end{solution}

\begin{problem}[A]
有 \(8\) 本互不相同的书，其中数学书 \(3\) 本，外文书 \(2\) 本，其它书 \(3\) 本. 若将这些书排成一列放在书架上，则数学书恰好排在一起，外文书也恰好排在一起的排法共有\fillinblank{}种（结果用数值表示）.
\end{problem}

\begin{answer}
\(1440\)
\end{answer}

\begin{solution}
把 \(3\) 本数学书看成一个整体，把 \(2\) 本外文书看成一个整体，再与其余 \(3\) 本书一起排列，共有 \(5!\) 种排列. 数学书这个整体内部有 \(3!\) 种排列，外文书这个整体内部有 \(2!\) 种排列. 因此排法总数为
\[
5!\cdot3!\cdot2!=120\cdot6\cdot2=1440
\]
\end{solution}

\begin{problem}[A]
把半径为 \(3\mathrm{~cm}\)，中心角为 \(\frac{2}{3}\pi\) 的扇形卷成一个圆锥形容器，这个容器的容积为\fillinblank{}\(\mathrm{~cm}^3\)（结果保留 \(\pi\)）.
\end{problem}

\begin{answer}
\(\frac{2\sqrt{2}}{3}\pi\)
\end{answer}

\begin{solution}
扇形半径是圆锥的母线长 \(3\mathrm{~cm}\). 扇形弧长为
\[
3\cdot\frac{2\pi}{3}=2\pi\mathrm{~cm}
\]
它卷成圆锥底面的圆周，所以 \(2\pi r=2\pi\)，得 \(r=1\mathrm{~cm}\). 圆锥高为
\[
h=\sqrt{3^2-1^2}=2\sqrt{2}\mathrm{~cm}
\]
因此容积为
\[
V=\frac{1}{3}\pi r^2h=\frac{2\sqrt{2}}{3}\pi\mathrm{~cm}^3
\]
\end{solution}

\setcounter{probnum}{14}
\begin{problem}[B]
已知 \(\bs{a}+\bs{b}=2\bs{i}-8\bs{j}\)，\(\bs{a}-\bs{b}=-8\bs{i}+16\bs{j}\). 那么 \(\bs{a}\cdot\bs{b}=\)\fillinblank{}.
\end{problem}

\begin{answer}
\(-63\)
\end{answer}

\begin{solution}
由
\[
\bs{a}=\frac{(\bs{a}+\bs{b})+(\bs{a}-\bs{b})}{2}
=\frac{-6\bs{i}+8\bs{j}}{2}
=-3\bs{i}+4\bs{j}
\]
以及
\[
\bs{b}=\frac{(\bs{a}+\bs{b})-(\bs{a}-\bs{b})}{2}
=\frac{10\bs{i}-24\bs{j}}{2}
=5\bs{i}-12\bs{j}
\]
所以
\[
\bs{a}\cdot\bs{b}=(-3)\cdot5+4\cdot(-12)=-15-48=-63
\]
\end{solution}

\begin{problem}[B]
\(\displaystyle\lim_{x\to 2}\left(\frac{4}{x^2-4}-\frac{1}{x-2}\right)=\)\fillinblank{}.
\end{problem}

\begin{answer}
\(-\frac{1}{4}\)
\end{answer}

\begin{solution}
当 \(x\ne2\) 时，
\[
\frac{4}{x^2-4}-\frac{1}{x-2}
=\frac{4-(x+2)}{(x-2)(x+2)}
=\frac{-(x-2)}{(x-2)(x+2)}
=-\frac{1}{x+2}
\]
因此
\[
\lim_{x\to2}\left(\frac{4}{x^2-4}-\frac{1}{x-2}\right)
=\lim_{x\to2}\left(-\frac{1}{x+2}\right)
=-\frac{1}{4}
\]
\end{solution}

\begin{problem}[B]
有 \(8\) 本互不相同的书，其中数学书 \(3\) 本，外文书 \(2\) 本，其它书 \(3\) 本. 若将这些书随机地排成一列放在书架上，则数学书恰好排在一起，外文书也恰好排在一起的概率为\fillinblank{}（结果用分数表示）.
\end{problem}

\begin{answer}
\(\frac{1}{28}\)
\end{answer}

\begin{solution}
将 \(3\) 本数学书捆成一个整体，将 \(2\) 本外文书捆成一个整体，再与其余 \(3\) 本书排列，满足条件的排法数为
\[
5!\cdot3!\cdot2!=1440
\]
全部 \(8\) 本互不相同的书随机排列，排法总数为 \(8!\). 因此所求概率为
\[
\frac{5!\cdot3!\cdot2!}{8!}
=\frac{1440}{40320}
=\frac{1}{28}
\]
\end{solution}

\begin{problem}[B]
某人有 \(10\) 万元，准备用于投资房地产或购买股票. 如果根据下面的盈利表进行决策：

\begin{center}
\begin{tabular}{|c|c|c|c|}
\hline
\multirow{2}{*}{自然状况} & 盈利（万元） & \multicolumn{2}{c|}{方案} \\
\cline{2-4}
 & 概率 & 购买股票 & 投资房地产 \\
\hline
巨大成功 & \(0.3\) & \(10\) & \(8\) \\
\hline
中等成功 & \(0.5\) & \(3\) & \(4\) \\
\hline
失败 & \(0.2\) & \(-5\) & \(-4\) \\
\hline
\end{tabular}
\end{center}

那么决策方案是\fillinblank{}.
\end{problem}

\begin{answer}
投资房地产
\end{answer}

\begin{solution}
购买股票的期望盈利为
\[
0.3\cdot10+0.5\cdot3+0.2\cdot(-5)=3.5
\]
万元；投资房地产的期望盈利为
\[
0.3\cdot8+0.5\cdot4+0.2\cdot(-4)=3.6
\]
万元. 按期望盈利较大的方案决策，应投资房地产.
\end{solution}

\section{解答题}

\begin{problem}
已知 \(\sin\left(\frac{\pi}{4}+\alpha\right)\sin\left(\frac{\pi}{4}-\alpha\right)=\frac{1}{6}\)，（\(\alpha\in\left(\frac{\pi}{2},\pi\right)\)），求 \(\sin 4\alpha\) 的值.
\end{problem}

\begin{solution}
\(\sin\left(\frac{\pi}{4}-\alpha\right)=\cos\left(\frac{\pi}{4}+\alpha\right)\)，所以
\[
\sin\left(\frac{\pi}{4}+\alpha\right)\sin\left(\frac{\pi}{4}-\alpha\right)
=\sin\left(\frac{\pi}{4}+\alpha\right)\cos\left(\frac{\pi}{4}+\alpha\right)
=\frac{1}{2}\sin\left(\frac{\pi}{2}+2\alpha\right)
=\frac{1}{2}\cos2\alpha
\]
由题设得 \(\cos2\alpha=\frac{1}{3}\). 又因为 \(\alpha\in\left(\frac{\pi}{2},\pi\right)\)，所以 \(2\alpha\in(\pi,2\pi)\)，从而
\[
\sin2\alpha=-\sqrt{1-\cos^22\alpha}
=-\sqrt{1-\frac{1}{9}}
=-\frac{2\sqrt{2}}{3}
\]
因此
\[
\sin4\alpha=2\sin2\alpha\cos2\alpha
=2\cdot\left(-\frac{2\sqrt{2}}{3}\right)\cdot\frac{1}{3}
=-\frac{4\sqrt{2}}{9}
\]
\end{solution}

\begin{problem}
在如图所示的直角坐标系中，一运动物体经过点 \(A(0,9)\)，其轨迹方程为 \(y=ax^2+c\)（\(a<0\)），\(D=(6,7)\) 为 \(x\) 轴上的给定区间.

\begin{enumerate}
\item 为使物体落在 \(D\) 内，求 \(a\) 的取值范围；
\item 若物体运动时又经过点 \(P(2,8.1)\)，问它能否落在 \(D\) 内？并说明理由.
\end{enumerate}

\bitmapfigure[max width=0.42\linewidth]{img/1996/shanghai_liberal/q20_fig1.png}
\end{problem}

\begin{solution}
\begin{enumerate}
\item 由点 \(A(0,9)\) 在轨迹上，得 \(c=9\)，所以轨迹方程为 \(y=ax^2+9\). 物体落到 \(x\) 轴上的位置满足
\(0=ax^2+9\)，\(x^2=-\frac{9}{a}\).
由于 \(a<0\)，取正半轴上的交点，其横坐标为 \(\sqrt{-\frac{9}{a}}\). “落在 \(D\) 内”表示该横坐标在 \((6,7)\) 内，因此
\[
6<\sqrt{-\frac{9}{a}}<7
\]
平方并注意 \(-a>0\)，得到
\(36<\frac{9}{-a}<49\)，\(\Longrightarrow\)，\(\frac{9}{49}<-a<\frac{1}{4}\).
故
\[
-\frac{1}{4}<a<-\frac{9}{49}
\]

\item 由于物体经过 \(P(2,8.1)\)，有
\[
8.1=4a+9
\]
解得 \(a=-\frac{9}{40}\). 又
\[
-\frac{1}{4}=-\frac{10}{40}<-\frac{9}{40}<-\frac{9}{49}
\]
所以该物体能落在 \(D\) 内.
\end{enumerate}
\end{solution}

\begin{problem}
如图，在二面角 \(\alpha-l-\beta\) 中，\(A\)，\(B\in\alpha\)，\(C\)，\(D\in l\)，\(ABCD\) 为矩形，\(P\in\beta\)，\(PA\perp\alpha\)，且 \(PA=AD\). \(M\)，\(N\) 依次是 \(AB\)，\(PC\) 的中点.

\begin{enumerate}
\item 求二面角 \(\alpha-l-\beta\) 的大小；
\item 求证：\(MN\perp AB\)；
\item 求异面直线 \(PA\) 与 \(MN\) 所成角的大小.
\end{enumerate}

\bitmapfigure[max width=0.42\linewidth]{img/1996/shanghai_liberal/q21_fig1.png}
\end{problem}

\begin{solution}
\begin{enumerate}
\item 连接 \(PD\). 因为 \(PA\perp\alpha\)，而 \(l\subset\alpha\)，所以 \(PA\perp l\)；又 \(AD\perp l\). 直线 \(PA\)，\(AD\) 相交且都垂直于 \(l\)，故平面 \(PAD\perp l\)，从而 \(PD\perp l\). 于是 \(AD\) 与 \(PD\) 分别位于平面 \(\alpha\)，\(\beta\) 内且都垂直于棱 \(l\)，所以 \(\angle PDA\) 是二面角 \(\alpha-l-\beta\) 的平面角. 在直角三角形 \(PAD\) 中，\(\angle PAD=90^\circ\)，且 \(PA=AD\)，所以 \(\angle PDA=45^\circ\). 二面角 \(\alpha-l-\beta\) 的大小为 \(45^\circ\).

\item 设 \(E\) 是 \(DC\) 的中点，连接 \(ME\)，\(NE\). 在矩形 \(ABCD\) 中，\(M\)，\(E\) 分别为 \(AB\)，\(DC\) 的中点，所以 \(ME\parallel AD\). 在三角形 \(PCD\) 中，\(N\)，\(E\) 分别为 \(PC\)，\(DC\) 的中点，所以 \(NE\parallel PD\). 由第（1）问知 \(AD\perp l\)，\(PD\perp l\)，故 \(ME\perp l\)，\(NE\perp l\). 直线 \(ME\)，\(NE\) 相交且都垂直于 \(l\)，所以 \(l\perp\) 平面 \(MNE\). 又 \(AB\parallel l\)，从而 \(AB\perp\) 平面 \(MNE\). 由于 \(MN\) 在平面 \(MNE\) 内，得到 \(MN\perp AB\).

\item 设 \(Q\) 是 \(DP\) 的中点，连接 \(NQ\)，\(AQ\). 在三角形 \(DPC\) 中，\(N\)，\(Q\) 分别为 \(PC\)，\(PD\) 的中点，因此
\(NQ\parallel DC\)，\(NQ=\frac{1}{2}DC\).
在矩形 \(ABCD\) 中，\(AM\parallel DC\)，且
\[
AM=\frac{1}{2}AB=\frac{1}{2}DC
\]
所以 \(QN\parallel AM\) 且 \(QN=AM\)，四边形 \(QNMA\) 是平行四边形，从而 \(AQ\parallel MN\). 因此异面直线 \(PA\) 与 \(MN\) 所成的角等于 \(PA\) 与 \(AQ\) 所成的锐角，即 \(\angle PAQ\). 三角形 \(PAD\) 是以 \(A\) 为直角顶点的等腰直角三角形，\(Q\) 是斜边 \(PD\) 的中点，所以中线 \(AQ\) 也是顶角平分线，\(\angle PAQ=45^\circ\). 故 \(PA\) 与 \(MN\) 所成角的大小为 \(45^\circ\).
\end{enumerate}
\end{solution}

\section{选做题：解答题}

\begin{problem}[A]
设 \(z=a+b\i\)（\(a\)，\(b\in\R\)，\(b\neq0\)，\(\i\) 为虚数单位），\(w=z+\frac{1}{z}\) 是实数，且 \(-1<w<2\).

\begin{enumerate}
\item 求 \(\left|z\right|\) 的值及 \(z\) 的实部的取值范围；
\item 设 \(u=\frac{1-z}{1+z}\). 求证：\(u\) 为纯虚数；
\item 求 \(w-u^2\) 的最小值.
\end{enumerate}
\end{problem}

\begin{solution}
\begin{enumerate}
\item 由
\[
\frac{1}{z}=\frac{a-b\i}{a^2+b^2}
\]
可得
\[
w=z+\frac{1}{z}
=\left(a+\frac{a}{a^2+b^2}\right)
+\left(b-\frac{b}{a^2+b^2}\right)\i
\]
因为 \(w\) 是实数，且 \(b\neq0\)，虚部必须为零，所以
\(b-\frac{b}{a^2+b^2}=0\)，\(\Longrightarrow\)，\(a^2+b^2=1\).
因此 \(|z|=1\)，并且 \(w=2a\). 由 \(-1<w<2\) 得
\[
-\frac{1}{2}<a<1
\]
所以 \(z\) 的实部的取值范围为 \(\left(-\frac{1}{2},1\right)\).

\item 由 \(b\neq0\) 和 \(a^2+b^2=1\) 可知 \(a\neq-1\)，所以 \(1+z\neq0\). 分母有理化得
\[
u=\frac{1-a-b\i}{1+a+b\i}
=\frac{(1-a-b\i)(1+a-b\i)}{(1+a+b\i)(1+a-b\i)}
=\frac{1-a^2-b^2-2b\i}{(1+a)^2+b^2}
\]
利用 \(a^2+b^2=1\) 以及
\((1+a)^2+b^2=2(a+1)\)，得到
\[
u=-\frac{b}{a+1}\i
\]
这里 \(-\frac{b}{a+1}\) 为实数，因此 \(u\) 为纯虚数.

\item 因为 \(u=-\frac{b}{a+1}\i\)，所以
\[
w-u^2=2a+\frac{b^2}{(a+1)^2}
=2a+\frac{1-a^2}{(a+1)^2}
=2a-\frac{a-1}{a+1}
=2\left(a+1+\frac{1}{a+1}\right)-3
\]
由 \(-\frac{1}{2}<a<1\)，令 \(t=a+1\)，则 \(t>0\). 由基本不等式
\[
t+\frac{1}{t}\geq2
\]
可得
\[
w-u^2\geq2\cdot2-3=1
\]
当且仅当 \(t=1\)，即 \(a=0\) 时取等号；此时可取 \(b=\pm1\)，且 \(w=0\) 满足题设范围. 因此 \(w-u^2\) 的最小值为 \(1\).
\end{enumerate}
\end{solution}

\reuseproblemnumber
\begin{problem}[B]
已知 \(A(-1,2)\) 为抛物线 \(C:y=2x^2\) 上的点，直线 \(l_1\) 过点 \(A\)，且与抛物线 \(C\) 相切. 直线 \(l_2:x=a\)（\(a\neq-1\)）交抛物线 \(C\) 于点 \(B\)，交直线 \(l_1\) 于点 \(D\).

\begin{enumerate}
\item 求直线 \(l_1\) 的方程；
\item 设 \(\triangle ABD\) 的面积为 \(S_1\)，求 \(|BD|\) 及 \(S_1\) 的值；
\item 设由抛物线 \(C\)，直线 \(l_1\)，\(l_2\) 所围成的图形的面积为 \(S_2\)，求证：\(S_1:S_2\) 的值为与 \(a\) 无关的常数.
\end{enumerate}
\end{problem}

\begin{solution}
\begin{enumerate}
\item 抛物线 \(C\) 的函数为 \(y=2x^2\)，所以
\(y'=4x\). 在点 \(A(-1,2)\) 处，切线斜率为 \(-4\). 因而
\[
y-2=-4(x+1)
\]
即
\[
4x+y+2=0
\]

\item 由 \(l_2:x=a\) 与抛物线 \(C\) 相交，得
\[
B=(a,2a^2)
\]
由 \(l_2:x=a\) 与 \(l_1\) 相交，得
\[
D=(a,-4a-2)
\]
所以 \(BD\) 是竖直线段，且
\[
|BD|=\left|2a^2-(-4a-2)\right|
=2(a+1)^2
\]
其中 \(a\neq-1\) 保证线段长度非零. 点 \(A\) 到直线 \(BD:x=a\) 的距离为
\(|a+1|\). 因此
\[
S_1=\frac{1}{2}|BD|\cdot|a+1|
=\frac{1}{2}\cdot2(a+1)^2|a+1|
=|a+1|^3
\]

\item 抛物线与切线的纵坐标差为
\[
2x^2-(-4x-2)=2(x+1)^2\geq0
\]
当 \(a>-1\) 时，所围区域在 \(x=-1\) 与 \(x=a\) 之间，所以
\[
S_2=\int_{-1}^{a}2(x+1)^2\,\mathrm{d}x
=\frac{2}{3}(a+1)^3
\]
当 \(a<-1\) 时，所围区域在 \(x=a\) 与 \(x=-1\) 之间，所以
\[
S_2=\int_{a}^{-1}2(x+1)^2\,\mathrm{d}x
=-\frac{2}{3}(a+1)^3
=\frac{2}{3}|a+1|^3
\]
综上，
\[
\frac{S_1}{S_2}
=\frac{|a+1|^3}{\frac{2}{3}|a+1|^3}
=\frac{3}{2}
\]
故 \(S_1:S_2=3:2\)，与 \(a\) 无关.
\end{enumerate}
\end{solution}

\section{解答题}

\begin{problem}
已知双曲线 \(S\) 的两条渐近线过坐标原点，且与以点 \(A(\sqrt{2},0)\) 为圆心，\(1\) 为半径的圆相切，双曲线 \(S\) 的一个顶点 \(A^\prime\) 与点 \(A\) 关于直线 \(y=x\) 对称.

\begin{enumerate}
\item 求双曲线 \(S\) 的方程；
\item 设直线 \(l\) 过点 \(A\)，斜率为 \(1\). 在双曲线 \(S\) 的上支上求点 \(B\)，使其与直线 \(l\) 的距离为 \(\sqrt{2}\)；
\item 求双曲线 \(S\) 上到直线 \(l:y=\frac{2\sqrt{5}}{5}\left(x-\sqrt{2}\right)\) 的距离为 \(\sqrt{2}\) 的点的个数.
\end{enumerate}

\bitmapfigure[max width=0.34\linewidth]{img/1996/shanghai_liberal/q22_fig1.png}
\end{problem}

\begin{solution}
\begin{enumerate}
\item 两条渐近线均过原点，可设其中一条为 \(y=kx\). 它与圆
\(\left(x-\sqrt{2}\right)^2+y^2=1\) 相切，所以圆心 \(A(\sqrt{2},0)\) 到直线 \(kx-y=0\) 的距离为 \(1\)，即
\[
\frac{|k|\sqrt{2}}{\sqrt{k^2+1}}=1
\]
平方得 \(2k^2=k^2+1\)，所以 \(k=\pm1\)，两条渐近线为 \(y=\pm x\). 因而 \(S\) 是等轴双曲线，且其一个顶点是点 \(A\) 关于 \(y=x\) 的对称点 \(A'(0,\sqrt{2})\). 双曲线的顶点在 \(y\) 轴上，所以可写为
\[
\frac{y^2}{a^2}-\frac{x^2}{a^2}=1
\]
由 \(a=\sqrt{2}\)，得到
\[
\frac{y^2}{2}-\frac{x^2}{2}=1
\]
即 \(y^2-x^2=2\).

\item 直线 \(l\) 过 \(A(\sqrt{2},0)\) 且斜率为 \(1\)，所以
\[
l:y=x-\sqrt{2}
\]
双曲线 \(S\) 的上支可表示为 \(y=\sqrt{x^2+2}\). 点 \(B\left(x,\sqrt{x^2+2}\right)\) 到 \(l\) 的距离为 \(\sqrt{2}\)，故
\[
\frac{\left|x-\sqrt{x^2+2}-\sqrt{2}\right|}{\sqrt{2}}=\sqrt{2}
\]
即 \(\left|x-\sqrt{x^2+2}-\sqrt{2}\right|=2\). 由于
\(\sqrt{x^2+2}>|x|\)，有 \(x-\sqrt{x^2+2}-\sqrt{2}<0\)，所以
\[
\sqrt{x^2+2}=x+2-\sqrt{2}
\]
两边平方得
\[
x^2+2=(x+2-\sqrt{2})^2
\]
即
\[
2(2-\sqrt{2})x=2-(2-\sqrt{2})^2=4(\sqrt{2}-1)
\]
从而 \(x=\sqrt{2}\). 此时右端 \(x+2-\sqrt{2}=2>0\)，且
\(y=\sqrt{x^2+2}=2\)，代回距离式可知原条件成立. 所以
\[
B=(\sqrt{2},2)
\]

\item 记 \(m=\frac{2\sqrt{5}}{5}\)，原直线为 \(l:y=mx-m\sqrt{2}=mx-\frac{2\sqrt{10}}{5}\). 与 \(l\) 平行且距离为 \(\sqrt{2}\) 的直线可写为 \(y=mx+c\)，并满足
\[
\frac{\left|c+\frac{2\sqrt{10}}{5}\right|}{\sqrt{m^2+1}}=\sqrt{2}
\]
由于 \(m^2+1=\frac{9}{5}\)，所以
\[
\left|c+\frac{2\sqrt{10}}{5}\right|=\frac{3\sqrt{10}}{5}
\]
从而
\(c_1=\frac{\sqrt{10}}{5}\)，\(c_2=-\sqrt{10}\).
记相应直线为
\(l_1:y=\frac{2\sqrt{5}}{5}x+\frac{\sqrt{10}}{5}\)，\(l_2:y=\frac{2\sqrt{5}}{5}x-\sqrt{10}\).
到 \(l\) 的距离为 \(\sqrt{2}\) 的双曲线上的点，恰为这两条平行线与双曲线的公共点. 将 \(l_1\) 代入 \(y^2-x^2=2\)，得
\[
x^2-4\sqrt{2}x+8=0
\]
其判别式为 \(32-32=0\)，所以 \(l_1\) 与双曲线有 \(1\) 个公共点. 将 \(l_2\) 代入，得
\[
x^2+20\sqrt{2}x-40=0
\]
其判别式为
\[
(20\sqrt{2})^2+160=960>0
\]
所以 \(l_2\) 与双曲线有 \(2\) 个公共点. 两条直线不同且平行，公共点不重合，故所求点的个数为 \(1+2=3\).
\end{enumerate}
\end{solution}

\begin{problem}
设 \(A_n\) 为数列 \(\{a_n\}\) 前 \(n\) 项的和，\(A_n=\frac{3}{2}(a_n-1)\)（\(n\in\N\)）. 数列 \(\{b_n\}\) 的通项公式为 \(b_n=4n+3\)（\(n\in\N\)）.

\begin{enumerate}
\item 求数列 \(\{a_n\}\) 的通项公式；
\item 若 \(d\in\{a_1,a_2,a_3,\cdots,a_n,\cdots\}\cap\{b_1,b_2,b_3,\cdots,b_n,\cdots\}\)，则称 \(d\) 为数列 \(\{a_n\}\) 与 \(\{b_n\}\) 的公共项. 将数列 \(\{a_n\}\) 与 \(\{b_n\}\) 的公共项，按它们在原数列中的先后顺序排成一个新的数列 \(\{d_n\}\)，证明数列 \(\{d_n\}\) 的通项公式为 \(d_n=3^{2n+1}\)（\(n\in\N\)）.
\end{enumerate}
\end{problem}

\begin{solution}
\begin{enumerate}
\item 当 \(n=1\) 时，\(A_1=a_1\)，所以
\[
a_1=\frac{3}{2}(a_1-1)
\]
解得 \(a_1=3\). 当 \(n\geq2\) 时，
\[
a_n=A_n-A_{n-1}
=\frac{3}{2}(a_n-1)-\frac{3}{2}(a_{n-1}-1)
=\frac{3}{2}(a_n-a_{n-1})
\]
整理得 \(a_n=3a_{n-1}\). 因此 \(\{a_n\}\) 是首项为 \(3\)、公比为 \(3\) 的等比数列，故
\(a_n=3^n\)，\((n\in\N)\).

\item
\begin{enumerate}
\item 对任意 \(n\in\N\)，\(a_{2n+1}=3^{2n+1}\)，所以 \(3^{2n+1}\) 是数列 \(\{a_n\}\) 的项. 又因为 \(3\equiv-1\pmod4\)，所以
\[
3^{2n+1}\equiv3\pmod4
\]
因此存在正整数 \(m\)，使得 \(3^{2n+1}=4m+3=b_m\). 于是 \(3^{2n+1}\) 是两个数列的公共项.

\item \(a_1=3\)，而 \(b_m=4m+3\geq7\)（\(m\in\N\)），所以 \(a_1\) 不是 \(\{b_n\}\) 的项. 对任意 \(n\in\N\)，
\[
a_{2n}=3^{2n}\equiv1\pmod4
\]
而 \(b_m\equiv3\pmod4\)，所以所有偶数项 \(a_{2n}\) 都不是 \(\{b_n\}\) 的项. 结合上一项，两个数列的公共项恰为
\(a_3\)，\(a_5\)，\(a_7\)，\(\ldots\)，\(a_{2n+1}\)，\(\ldots\).
按原数列先后顺序排列后，
\(d_n=a_{2n+1}=3^{2n+1}\)，\((n\in\N)\).
\end{enumerate}
\end{enumerate}
\end{solution}
